\chapter{Literature Review and Theoretical Gap}

The theory developed in this dissertation sits at the intersection of five
literatures: classical power-system stability, spectral and network
interpretations of synchronization, second-order and quadratic eigenvalue
analysis, nonlinear eigenvalue problems generated by state condensation, and
converter-dominated stability. The purpose of this chapter is not to catalogue
every contribution in those areas. It is to identify the exact conceptual gap
that motivates a damping-operator hierarchy.

\section{Systematic Review Contract}

In the expanded dissertation this chapter should become a systematic review,
not only a narrative survey. The review protocol should report search strings,
databases, dates, inclusion and exclusion criteria, preprint versus
peer-reviewed status, extraction fields, and a method-assumption-data-certificate
matrix. The central taxonomy is the operator hierarchy
\begin{equation}
  \begin{aligned}
  D=\gamma M
  &\rightarrow
  [\Lt,\Dt]=0
  \rightarrow
  D\ {\rm fixed\ and\ non\!-\!commuting}
  \rightarrow
  \Pi(s)\ {\rm rational} \\
  &\rightarrow
  G_T(s)=T(s)^{-1}
  \rightarrow
  \Lambda_\varepsilon^{\mathcal A}(T)
  \rightarrow
  {\rm nonlinear\ memory}.
  \end{aligned}
  \label{eq:lit_operator_taxonomy}
\end{equation}
The first arrow marks the move from uniform damping-to-inertia ratio to the
larger decouplable class. The second marks modal coupling within a constant
QEP. The third marks controller or auxiliary-state condensation. The fourth
marks the Green function that carries finite actions. The fifth marks
structured fragility under physical perturbation families. The last arrow
marks the move from linearized spectral memory to genuinely nonlinear memory.
This taxonomy is the review's organizing principle.

\section{Classical Small-Signal Stability}

Classical power-system stability theory starts from an operating point and
studies the linearized differential-algebraic model around that point. The
standard engineering workflow is well established: solve the power flow,
initialize dynamic models, linearize, compute eigenvalues, inspect damping
ratios, and use eigenvectors or participation factors to interpret the critical
modes \cite{kundur1994power,rogers2000oscillations}. The IEEE/CIGRE stability
definition paper formalized the classification of rotor-angle, frequency, and
voltage stability and placed small-signal stability within a broader taxonomy
of disturbances, time scales, and mechanisms \cite{kundur2004definition}.

For the present thesis, the central fact is that classical modal analysis is
already an operator theory. A mode is not merely an oscillation observed in a
time-domain trace. It is a spectral object of a linearized system. Its damping
ratio is computed from the real and imaginary parts of a pole. Its engineering
interpretation depends on how the corresponding eigenvectors project onto
angles, speeds, exciters, governors, loads, and controls. Thus, any reduced
model that claims to preserve damping must preserve the relevant spectral
mechanism, not just the visual shape of a response.

\section{Graph and Synchronization Perspectives}

Power networks are also graphs. Under lossless or nearly lossless assumptions,
the active-power Jacobian has a Laplacian structure: off-diagonal terms encode
coupling between connected buses and diagonal terms enforce row-sum balance.
This structure links power-system dynamics to synchronization theory. The work
of Doerfler, Chertkov, and Bullo gives a concise synchronization condition for
complex oscillator networks and smart grids, clarifying how coupling strength,
phase cohesiveness, and network structure interact \cite{dorfler2013synchronization}.

The graph viewpoint is indispensable because it turns a high-dimensional
network into interpretable objects: algebraic connectivity, effective
resistance, cut structure, mode shapes, Kron-reduced Laplacians, and weighted
network strength. Algebraic connectivity traces back to Fiedler's second
Laplacian eigenvalue and its relation to graph connectivity
\cite{fiedler1973algebraic}. Effective-resistance design provides another
view of robustness by penalizing large pairwise electrical distances and by
linking network reinforcement to convex edge-weight allocation
\cite{ghosh2008effective}. Spectral sparsification shows why effective
resistance is also a structural sampling quantity: preserving selected
quadratic forms preserves the Laplacian action that matters to many slow-mode
metrics \cite{spielman2011sparsification}. Kron reduction gives the matching
electrical-network operation for eliminating passive buses through a Schur
complement while preserving a reduced Laplacian
\cite{dorfler2013kron}.

These tools motivate planning intuitions: reinforce a weak line, improve a
weak cut, place a grid-forming converter at a strategic bus, or raise a
critical graph eigenvalue. Those intuitions are valuable. The problem is that
damping is not always a graph eigenvalue. Damping may be local, heterogeneous,
non-proportional, controller-mediated, or frequency dependent. The graph
literature gives the stiffness skeleton; the present thesis asks which damping
operator must be attached to that skeleton.

This boundary is important for the originality claim. For example, Guo, Zhao,
and Low already connect the scaled graph Laplacian spectrum with primary
frequency regulation under uniform damping-to-inertia ratios
\cite{guo2018graph}. The present thesis uses that kind of result as a known
floor. Its contribution is not the existence of Laplacian modes, but the
operator hierarchy that decides when those modes remain damping coordinates,
when rough damping-to-inertia fields mix them, and when controller self-energy
must replace a static graph screen.

\section{Beyond Scalar Laplacians: Cellular Sheaves}

An important structural extension is to look beyond scalar graph
Laplacians. The closest mature mathematical object is the cellular sheaf
Laplacian. Hansen and Ghrist develop spectral sheaf theory by lifting the
ordinary graph Laplacian to a Hodge Laplacian acting on a sheaf of vector
spaces over a cell complex \cite{hansen2019sheaf}. In that setting, vertices
and edges may carry vector spaces, and restriction maps encode how local data
are compared across cells. The theory connects spectral data to sheaf
cohomology and includes analogues of interlacing, sparsification, effective
resistance, and synchronization.

This dissertation does not claim a completed sheaf theory of converter-rich
power grids. The point is more precise: a scalar Laplacian assigns one
number-valued state to each bus, whereas a converter-rich bus carries internal
controller coordinates. A sheaf model is a natural candidate for a future
operator in which the bus stalk contains retained electrical variables,
converter-control states, or reduced controller coordinates, and each line
restriction map describes which combinations are electrically compatible. Such
an operator could turn the present scalar mode-family labels and
self-energy-dominance diagnostics into a structured vector-bundle-like object.

The risk is equally clear. A sheaf Laplacian is useful only if it detects a
weak node, weak link, or hidden margin that the scalar \(L\), the fixed-\(D\)
QEP, and the Schur self-energy screen do not already detect. In this thesis it
is therefore placed as a speculative theoretical extension, not as evidence
for any current planning claim.

\section{Quadratic Eigenvalue Problems and Second-Order Models}

The swing-equation model is naturally second order. In compact form it leads to
a quadratic eigenvalue problem,
\begin{equation}
  P(s)v=(s^2M+sD+L)v=0.
  \label{eq:lit_qep}
\end{equation}
Tisseur and Meerbergen's survey establishes the QEP as a mature mathematical
object, with well-developed linearizations, spectral properties, perturbation
formulas, and numerical methods \cite{tisseur2001qep}. This literature matters
because the power-system damping problem should not be reduced too quickly to a
first-order matrix. The second-order structure carries physical information:
mass or inertia, damping, and stiffness are distinct operators.

The distinction between proportional and non-proportional damping is also
well known in structural dynamics. If damping and stiffness share a modal
basis, modal decoupling is exact. If they do not, modal analysis becomes
coupled, and non-proportionality indices are needed to measure the error of
decoupled approximations \cite{conde2021nonproportionality}. This thesis uses
the same mathematical distinction in a power-system setting. The first
departure from graph-frequency intuition is not necessarily converter
dynamics; it can already occur when \(D\) is fixed but heterogeneous relative
to \(M\).

\section{Nonlinear Eigenvalue Problems and Schur Condensation}

The next mathematical layer is the nonlinear eigenvalue problem (NEP). Guettel
and Tisseur survey matrix-valued functions \(T(s)\) whose dependence on the
spectral parameter is nonlinear and show that such problems arise naturally
when delayed, damped, or condensed dynamics are represented without expanding
all states \cite{guettel2017nep}. A Schur complement is a basic route from a
larger linear eigenvalue problem to a smaller nonlinear one. If the eliminated
block is \(A_{cc}\), the retained operator contains the resolvent
\((sI-A_{cc})^{-1}\).

For computation, contour integral methods and rational-interpolation
viewpoints are particularly relevant. Brennan, Embree, and Gugercin recast
contour integral methods for NEPs in a systems-theoretic framework using
Loewner matrix pencils, allowing interpolation away from infinity and linking
eigenvalue extraction to system realization and modal truncation
\cite{brennan2023contour}. This is important for a Schur self-energy thesis
because \(\Pi(s)\) and its damping-channel projection \(\Sigma(s)\) are
transfer-like rational objects. If the
controller block is not fully known, a Loewner or data-driven contour
formulation is a plausible route for identifying the visible part of
\(\Pi(s)\) from frequency-domain samples rather than from explicit
state-space matrices.

Recent power-system work also strengthens the numerical side of the proposed
planning workflow. Bouterakos, McKeon, and Tzounas develop continuation-based
eigenvalue tracking for large-scale power-system models, supporting sparse
representations and both explicit and semi-implicit DAE forms
\cite{bouterakos2025tracking}. That result is directly aligned with the local
pole-motion map used later in this dissertation: a planning screen that changes
penetration, line weights, inertia, damping, or controller gains must track
selected eigenvalue trajectories rather than repeatedly rediscovering the entire
spectrum from scratch.

This observation is the mathematical hinge of the dissertation. Converter and
auxiliary controller states can be retained explicitly, producing a large
first-order model. They can also be condensed, producing a smaller nonlinear
operator. The condensation is exact away from poles of the eliminated block,
but the price of exactness is frequency dependence. A constant damping matrix
cannot store the poles, zeros, and residues introduced by the eliminated
controller states. The reduced controller object is therefore a rational
self-energy; its damping-channel projection is the rational object denoted
\(\Sigma(s)\).

\section{Pseudospectra, Fragility, and Model Adequacy}

Nonlinear eigenvalue theory also raises a robustness question. A nominal root
set may look safe while nearby structured perturbations can move a pole to a
protected boundary, especially when the retained operator is non-normal. The
classical pseudospectral viewpoint already warns that eigenvalues alone can
miss transient amplification and perturbation sensitivity
\cite{trefethen2005spectra}. Recent work on nonlinear spectral problems
extends the computational discussion of spectra and pseudospectra to broad NEP
settings and emphasizes spectral pollution and invisibility risks in
approximations \cite{colbrook2025universal}.

For this dissertation, the distinction is that power-system perturbations must
be structured. A generic matrix perturbation is not the same as uncertainty in
line susceptance, PLL gains, converter delays, virtual inertia, damping, or an
identified self-energy. This motivates a physical-action pseudospectrum
\(\Lambda_\varepsilon^{\mathcal A}(T)\) and a fragility radius: the smallest
admissible perturbation that drives a retained root to a protected boundary.
The related model-adequacy problem is contour based. If a reduced operator
\(\widehat T\) is used for a decision, it should preserve the protected
eigenvalue count on a contour before its derivative signs are trusted.

\section{Sign Characteristics and Krein-Type Mechanisms}

The mathematical language closest to the proposed controller-resolvent
collision is not a new, free-floating ``Krein signature for NEPs.'' The citable
foundation is the sign characteristic of Hermitian matrix polynomials and
self-adjoint meromorphic matrix functions. Gohberg, Lancaster, and Rodman give
the classical matrix-polynomial framework, including sign characteristics of
spectral data \cite{gohberg1982matrixpolynomials}. Their related work on
self-adjoint meromorphic matrix functions is closer to the rational structure
created by Schur condensation \cite{gohberg1983meromorphic}. Mehrmann,
Noferini, Tisseur, and Xu develop sign characteristics for Hermitian matrix
polynomials and clarify the role of structure-preserving transformations
\cite{mehrmann2016sign}. Kollár and Miller's Evans--Krein survey provides a
systems-oriented view in which rational Krein functions track spectral
collisions \cite{kollar2014graphical}.

Those references do not already solve the converter problem. They mostly live
in Hamiltonian, self-adjoint, or Hermitian-polynomial settings. Converter
self-energies are rational, dissipative, and generally non-normal. The thesis
therefore treats sign characteristic as a foundation to be adapted, not as a
ready-made theorem. The open bridge is to show when a rational
controller-condensed operator has an operational sign collision that coincides
with weak-grid PLL or controller-resonance instability.

\section{Converter-Dominated Stability}

The increased penetration of converter-interfaced generation motivates a
reconsideration of stability definitions and modeling practice. The revisited
IEEE/CIGRE classification explicitly adds converter-driven and electric
resonance phenomena to the stability vocabulary
\cite{hatziargyriou2021stability}. IEEE Std 2800-2022 reflects the same
industrial shift by specifying performance and interconnection requirements for
IBRs connected to transmission systems \cite{ieee2800_2022}.

Within the research literature, grid-following and grid-forming controls are
not interchangeable dynamic objects. Grid-following converters commonly depend
on a PLL for synchronization, so PLL dynamics can contribute to small-signal
instability under weak-grid conditions. Yang, Huang, Xin, and Ju analyze how
placing grid-forming converters can improve the small-signal stability of
PLL-integrated systems by modifying grid strength and Laplacian-related
quantities \cite{yang2021gfm}. Comparative studies of grid-forming and
grid-following inverters further show that control architecture changes the
state-space model and its eigenvalues \cite{ding2021comparative}.

Inertia and primary-control placement studies sharpen the same point. Optimal
virtual inertia allocation can be formulated as an \(H_2\)-type network
performance problem \cite{poolla2017virtual}; disturbance propagation depends
strongly on where inertia is located relative to slow network modes
\cite{pagnier2019inertia_location}; and perturbation theory shows that inertia
and primary-control placement are not interchangeable under heterogeneous
parameters \cite{pagnier2019optimal}. These results support the thesis
distinction between adding inertia, adding damping, and modifying the graph.

Passivity and port-Hamiltonian viewpoints provide a second, complementary
control language. Passivity-based grid-forming work formulates converter
controllers so that local subsystems have dissipativity properties that can be
interconnected modularly. He and D{\"o}rfler derive decentralized stability
conditions for grid-forming converters using passivity indices
\cite{he2024passivity}. Huang et al. combine small-gain and small-phase
arguments to obtain decentralized stability checks for heterogeneous
power-electronics-dominated systems \cite{huang2024gainphase}, while Gui,
Subedi, and Xue use a port-controlled Hamiltonian form for passivity-based
grid-forming control of distributed energy resources \cite{gui2024passivity}.
Recent matrix-valued passivity indices make the same issue sharper for MIMO
ports: scalar indices can discard the direction and coupling structure of a
passivity deficit \cite{ru2026matrix}. Dynamic passivity multipliers further
show why frequency-dependent certificates can be less conservative than strict
static passivity tests for converter-dominated grids \cite{gorbunov2026dynamic}.
In the notation of this thesis, this literature says that the Hermitian,
weighted directional, or multiplier-based behavior of the condensed
self-energy \(\Sigma(s)\) should not be a decorative diagnostic. It is a
candidate certificate for controller retuning, provided that the passivity
screen is tied back to the retained pole motion and to full-model validation.

These studies justify the engineering relevance of the present theory, but
they do not remove the need for a general operator language. A placement rule
derived for a particular control model or grid-strength metric may not
transfer to a different controller, parameter regime, or retained-state
partition. The common mathematical issue is not the label GFL or GFM by
itself; it is whether the reduced model preserves the spectral contribution of
the eliminated control states.

\section{Hosting Capacity, Network Reinforcement, and Braess Effects}

Hosting capacity is usually introduced as a planning quantity: how much
distributed generation, load, or converter-interfaced equipment can be added
before a registered performance limit is violated. Bollen and Roennberg frame
hosting capacity as a distribution-planning tool under uncertainty
\cite{bollen2017hosting}. Later dynamic hosting-capacity studies emphasize
that snapshot limits can miss the time duration, severity, and controller
response of violations \cite{oliveira2019dynamicHosting,jain2020dynamicHosting}.
Those works motivate the same methodological principle used in this thesis:
capacity is not a single number until the limiting metric has been named.

The present dissertation uses the phrase dynamic hosting capacity in a
different but compatible sense. Instead of asking only how long a voltage or
thermal violation lasts, it asks how much IBR conversion can be admitted before
a small-signal spectral margin is lost. Static hosting constraints remain
necessary, but they are insufficient in low-inertia systems because a
voltage-safe operating point can still have an unacceptable damping ratio,
resolvent gain, or controller-resonance exposure.

Recent stability-manifold work for converter-dominated systems reinforces the
need to map controller-parameter regions under multiple operating scenarios
\cite{conte2026stabilitymanifolds}. The distinction here is mechanism: the
present thesis seeks not only a sampled stability boundary, but an analytic
normal, a physical self-energy label, an error certificate, and a minimum
admissible action associated with that boundary.

Network reinforcement has a similar ambiguity. In a purely stiffness-oriented
graph view, increasing a line weight or adding a line seems beneficial.
Braess' original traffic result already showed that adding capacity can worsen
system performance under decentralized routing \cite{braess2005traffic}.
Power-system analogues have been studied for oscillator synchronization and
electric grids. Witthaut and Timme showed that adding links in oscillator
networks can destroy synchronization \cite{witthaut2012braess}, while Coletta
and Jacquod analyzed how added couplings can enhance or reduce linear
stability in coupled-oscillator and power-grid models \cite{coletta2016braess}.
Schaefer and coauthors further demonstrated Braess' paradox in AC power-grid
experiments and large-grid studies, emphasizing that line upgrades can
increase loading elsewhere in the network \cite{schafer2022braess}.

For this thesis, the Braess literature gives a warning rather than a license
to overclaim. A line reinforcement that increases modal frequency but reduces
damping ratio is a Braess-like damping-margin reversal, not automatically a
classical flow Braess paradox. A controller-condensed case in which the line
move pulls a network-family pole toward a nearby control pole is a stronger
mechanistic claim, but it still requires a registered NEP or full-model
validation gate. This calibrated vocabulary is used in the planning chapters.

\section{Computational Modeling and Reproducible Dynamics}

Detailed simulation environments are essential for testing the theory. ANDES,
for example, implements a hybrid symbolic-numeric framework for power-system
DAE modeling and analysis, including eigenvalue analysis and time-domain
simulation \cite{cui2021andes}. Such tools make it possible to build
converter-rich cases, extract Jacobians, compare full and reduced spectra, and
test whether a proposed self-energy approximation is valid.

The next modeling frontier is to identify the port self-energy when the
internal converter model is unavailable. Recent gray-box converter
identification work estimates voltage-source converter parameters from terminal
time-series data and validates the resulting model with frequency-domain
metrics \cite{darii2026graybox}. Structure-preserving identification of
port-Hamiltonian DAEs points in a complementary direction: learn from noisy
input-output data while preserving interconnection and passivity properties
\cite{hagelaars2026identification}. These lines support the thesis requirement
that an identified \(\widehat\Pi(s)\) must carry an uncertainty bound, not just
a fitted curve.

This dissertation, however, intentionally separates theory from validation.
The present manuscript develops the operator framework. A later empirical
chapter or companion article may use ANDES or another DAE platform to test the
framework on specific systems. That separation is methodological: without a
clear theoretical object, simulation can produce plausible plots without
explaining which damping mechanism has been preserved.

\section{Theoretical Gap}

The reviewed literatures leave a precise gap. Classical small-signal stability
provides modal damping ratios and participation factors. Spectral graph theory
provides an interpretable stiffness skeleton. QEP theory explains
second-order fixed-coefficient dynamics. NEP theory explains frequency-dependent
operators created by condensation. Converter-dominated stability shows why
controller dynamics matter. What is missing is a unified damping-operator
hierarchy that tells the engineer which level is mathematically sufficient for
a given claim.

The gap can be stated as four questions.

\begin{enumerate}
\item When does a graph mode carry a true damping certificate rather than only
a frequency screen?
\item When does fixed non-commuting damping create enough modal coupling to
invalidate decoupled modal reasoning?
\item When do controller states require a rational self-energy \(\Pi(s)\),
and when can its damping-channel projection \(\Sigma(s)\) replace a constant
reduced damping matrix?
\item When does a dynamic Green function or structured fragility radius change
the decision that would be made from a scalar graph or damping metric?
\end{enumerate}

The remainder of the dissertation answers these questions by developing the
hierarchy from uniform damping \(D=\gamma M\), to the broader commuting
constant-\(D\) class, to fixed non-commuting damping, to rational
self-energy \(\Pi(s)\), to Green-function action calculus, to structured
fragility, and finally to nonlinear memory.
